\documentclass[11pt,a4paper]{article}

\usepackage[utf8]{inputenc}
\usepackage{amsmath,amssymb,amsthm}
\usepackage{mathtools}
\usepackage[margin=1in]{geometry}
\usepackage{enumitem}
\usepackage{hyperref}

% Operators
\DeclareMathOperator{\Exp}{Exp}
\DeclareMathOperator{\Log}{Log}
\DeclareMathOperator{\Retr}{Retr}
\DeclareMathOperator{\Image}{Im}
\DeclareMathOperator{\Vol}{Vol}
\DeclareMathOperator{\tr}{tr}
\DeclareMathOperator{\SE}{SE}
\DeclareMathOperator{\SO}{SO}
\DeclareMathOperator{\so}{\mathfrak{so}}
\DeclareMathOperator{\se}{\mathfrak{se}}
\DeclareMathOperator{\Ad}{Ad}
\DeclareMathOperator{\ad}{ad}
\DeclareMathOperator{\KL}{KL}

% Commands
\newcommand{\R}{\mathbb{R}}
\newcommand{\E}{\mathbb{E}}
\newcommand{\N}{\mathcal{N}}
\newcommand{\D}{\mathcal{D}}
\newcommand{\Loss}{\mathcal{L}}
\newcommand{\Mphi}{\mathcal{M}_{\phi}}
\newcommand{\Mset}{\mathcal{M}}
\newcommand{\Mpose}{\mathcal{M}_\phi^{\text{pose}}}
\newcommand{\Mbpose}{\mathcal{M}_\phi^{\text{b-pose}}}
\newcommand{\Tphi}{T_\phi}
\newcommand{\Hpose}{H_\phi^{\text{pose}}}
\newcommand{\Hbpose}{H_\phi^{\text{b-pose}}}
\newcommand{\Jpose}{J_{\text{pose}}}
\newcommand{\JQ}{J^Q}
\newcommand{\GQ}{G_Q}
\newcommand{\GQbar}{\bar G_Q}
\newcommand{\psimap}{\psi}
\newcommand{\Dpsi}{D_\psimap}

\title{Core Mathematical Formulation\\
\large Pose-Extended SMCDP v5.1: $\SE(3)$ Self-Model on Bounded Chart with OU Forward\\
\normalsize (Corrected: Euclidean score convention, Convention-1 reverse SDE, honest feasibility tiering)}
\author{}
\date{}

\begin{document}

\maketitle

\paragraph{Document version note.} This is v5.1, a corrected revision of v5 addressing five mathematical and three narrative issues identified in review. The substantive design (bounded chart + OU forward + IK-free reference) is unchanged from v5; the corrections concern convention precision, claim honesty, and notational consistency.

\paragraph{Corrections from v5.}
\begin{enumerate}[noitemsep, leftmargin=*]
    \item \textbf{Finite-$K$ reference mismatch made explicit.} v5 occasionally claimed ``exact forward-reference match.'' Corrected: the OU process admits a closed-form stationary $p_\infty = \N(0, \GQbar^{-1})$ — this is the structural property absent in v4/v5-option-X. The forward marginal $p_K$ approaches $p_\infty$ only as $\tau(K) \to \infty$; at finite $K$ the mismatch is controlled by $\alpha(K) = e^{-\tau(K)/2}$ and must be measured empirically (\S\ref{sec:forward-reference-mismatch}).
    \item \textbf{Arithmetic for $\beta_f$-$\alpha$ relation corrected.} v5 erroneously paired $\beta_f = 10$ with $\alpha(K) \approx 0.007$. Correct: $\beta_f = 10 \Rightarrow \tau(K) = 5 \Rightarrow \alpha(K) = e^{-2.5} \approx 0.082$. The value $\alpha \approx 0.007$ requires $\beta_f \approx 20$. (\S\ref{sec:beta-schedule}, table updated.)
    \item \textbf{Score convention fixed to Euclidean.} v5 wrote the reverse SDE with $s_\theta^u$ preceded by $\GQbar^{-1}$ (suggesting Euclidean score) but wrote the DSM target with an extra $\GQ^{-1}$ factor (suggesting Riemannian gradient). v5.1 fixes: $s_\theta^u := \nabla_u \log p_r$ (Euclidean), and all target/reverse expressions are consistent with this definition (\S\ref{sec:score-conventions}).
    \item \textbf{Reverse SDE sign convention pinned to Convention 1.} Forward time $r$, $\Delta r < 0$ on reverse pass. Standard Score-SDE style. v5 implicitly mixed Convention 1 and 2; v5.1 is unambiguous (\S\ref{sec:reverse-sde}).
    \item \textbf{``Geodesic random walk'' removed.} v5 used this term loosely. The sampler is a chart-space OU sampler with graph retraction; not an intrinsic Riemannian Brownian motion under $\GQ$. v5.1 uses precise terminology (\S\ref{sec:sampling}).
    \item \textbf{Pose-block term reclassified as auxiliary regularizer.} The $\Log_{\SE(3)}$ component is not an exact forward-transition score (since $T = \widetilde T_\phi(u, z_e)$ is a deterministic function of $u$). v5.1 names it ``pose-geometric consistency regularizer'' rather than ``hybrid DSM target.'' (\S\ref{sec:loss}).
    \item \textbf{Conditioning channel claim corrected.} v5 said ``guidance is the sole conditioning channel.'' This is false: the score network input already contains $(T_\text{start}, T_\text{target}, z_e)$. v5.1: guidance is the \emph{non-network, non-IK explicit conditioning mechanism}, not the sole channel (\S\ref{sec:guidance}).
    \item \textbf{Feasibility claim tiered.} v5 claimed ``feasibility by construction.'' v5.1 separates this into (a) joint feasibility (true by chart structure), (b) learned-manifold adherence (true by retraction, w.r.t.\ $\Tphi$ not true robot), and (c) true physical feasibility (depends on self-model accuracy; collision/torque/dynamics not addressed) (\S\ref{sec:feasibility-tiers}).
    \item \textbf{``Riemannian SDE on the manifold'' claim weakened.} The forward SDE is a constant-coefficient OU on the chart $u$; it is not a metric-aware Riemannian Brownian motion under $\GQ(u, z_e)$. $\GQ$ enters only through weighting (loss, guidance preconditioning, tangent identity). v5.1 makes this distinction explicit (\S\ref{sec:forward-sde-positioning}).
    \item \textbf{Reference-marginal match metric promoted to required evaluation.} $\KL(p_K \| p_\text{ref})$ or $W_2$ must be reported; it justifies the IK-free reference choice for the chosen $\beta_f$ (\S\ref{sec:eval}).
\end{enumerate}

\paragraph{What v5.1 retains from v5 (the core contribution).}
\begin{itemize}[noitemsep]
    \item Bounded chart $u = \tanh^{-1}(\ldots)$ giving joint feasibility by construction (\S\ref{sec:bounded-chart}).
    \item OU forward with closed-form Gaussian transition kernel (\S\ref{sec:forward-sde}).
    \item Closed-form stationary $p_\infty = \N(0, \GQbar^{-1})$ as the structural basis for IK-free reference, with $\GQbar$ data-independent and conditioning-independent (\S\ref{sec:reference-metric}, \S\ref{sec:reference-distribution}).
    \item Closed-form OU score as the chart-block training target (exact, not Varadhan-approximated) (\S\ref{sec:loss}).
    \item Self-model graph manifold $\Mbpose(z_e) = \{(u, T) : T = \Tphi(\psi(u), z_e)\}$ preserved (\S\ref{sec:manifold}).
    \item Multi-component reward guidance for trajectory quality (\S\ref{sec:guidance}).
\end{itemize}

\paragraph{Notation.} Diffusion time $r \in [0, K]$, forward direction. Reverse pass uses $\Delta r < 0$ (Convention 1). Chart variable $u \in \R^{n_q}$ throughout; physical joint $q = \psi(u) \in (q_\text{min}, q_\text{max})$ is the lifted readout. Constant reference metric $\GQbar$ versus position/embodiment-dependent induced metric $\GQ(u, z_e)$; these are distinct objects with distinct roles.

\section{Bounded chart parameterization}
\label{sec:bounded-chart}

\paragraph{Chart map.}
\begin{equation}
q_i = \psi_i(u_i) = q_{\text{mid},i} + \frac{q_{\text{range},i}}{2} \tanh(u_i), \quad u_i \in \R.
\end{equation}
$\psi : \R^{n_q} \to (q_\text{min}, q_\text{max})$ is a smooth diffeomorphism onto the open interior. The boundary is approached only as $|u_i| \to \infty$.

\paragraph{Chart differential.}
\begin{equation}
\Dpsi(u) = \mathrm{diag}\bigl((q_{\text{range},i}/2)\, \mathrm{sech}^2(u_i)\bigr) \succ 0 \quad \forall u.
\end{equation}

\paragraph{Joint feasibility by construction.} For finite $u$, $\psi(u) \in (q_\text{min}, q_\text{max})$ strictly. This is independent of training data, noise scale, or extrapolation.

\paragraph{Saturation regime.} For $|u_i| > 3$, $\mathrm{sech}^2(u_i) < 0.01$. $\Dpsi$ near-singular; chart-space dynamics ill-conditioned. Practical operation should keep $\|u\|_\infty \lesssim 2$. The chart enforces the boundary as a hard limit, not as the operating regime.

\section{Self-model learning}
\label{sec:self-model}

Unchanged from v4. The self-model is defined on physical $q$:
\begin{equation}
\Tphi(q, z_e) = T_\text{analytic}(q, z_e) \cdot \exp_{\SE(3)}\!\bigl(\xi_\phi(q, z_e)^\wedge\bigr).
\end{equation}
Composition: $\widetilde T_\phi(u, z_e) := \Tphi(\psi(u), z_e)$.

Constraint function: $g_\phi(x, z_e) = \Log_{\SE(3)}(\Tphi(q, z_e)^{-1} T)$.

\section{Manifold on bounded chart}
\label{sec:manifold}

\begin{equation}
\Mbpose(z_e) = \bigl\{(u, T) : T = \widetilde T_\phi(u, z_e)\bigr\}.
\end{equation}
Embedding: $\Hbpose(u, z_e) = (u, \widetilde T_\phi(u, z_e))$. The first ambient block is the chart $u$ (not $\psi(u)$).

\section{Differential structure}
\label{sec:tangent}

\paragraph{Chart-coordinate body-frame Jacobian.}
\begin{equation}
\JQ(u, z_e) = \Jpose(\psi(u), z_e) \cdot \Dpsi(u) \in \R^{6 \times n_q}.
\end{equation}

\paragraph{Tangent map.}
\begin{equation}
J_H^{\text{b-pose}}(u, z_e) = \begin{bmatrix} I_{n_q} \\ \JQ(u, z_e) \end{bmatrix}.
\end{equation}

\paragraph{On-manifold tangent identity.}
\begin{equation}
J_g \cdot J_H^{\text{b-pose}} = 0 \quad \text{on } \Mbpose.
\end{equation}

\section{Induced metric (Choice A)}
\label{sec:metric}

\paragraph{Pose-tangent inner product.}
\begin{equation}
W = \mathrm{diag}(W_p I_3, W_R I_3), \quad W_p = \sigma_p^{-2} = 400, \quad W_R = \sigma_R^{-2} = 100.
\end{equation}

\paragraph{Induced chart-coordinate metric.}
\begin{equation}
\GQ(u, z_e) = I_{n_q} + {\JQ}^\top W \JQ \succeq I_{n_q}.
\end{equation}

\paragraph{Role of $\GQ$ in v5.1.} $\GQ$ enters in three places, none of them as the diffusion noise covariance:
\begin{enumerate}[noitemsep]
    \item \textbf{Loss weighting}: the score-matching loss can be $\GQ$-weighted (or weighted by $\GQ^{-1}$ or $I$) — ablation choice.
    \item \textbf{Guidance preconditioning}: $G_\text{traj}^{-1} \nabla_u R$ as a natural-gradient-style guidance (this is heuristic, not from exact reverse SDE).
    \item \textbf{Tangent verification}: $J_g \cdot J_H^{\text{b-pose}} = 0$ identity uses the structure of $\GQ = I + {\JQ}^\top W \JQ$.
\end{enumerate}
$\GQ$ \emph{does not} enter the forward SDE noise covariance — that role is played by the constant $\GQbar$.

\subsection{Adaptive Tikhonov}
\label{sec:tikhonov}
$\GQ^\text{reg} = \GQ + \lambda(u, z_e) I$, $\lambda = c_\lambda \tr(\GQ)/n_q$, $c_\lambda \approx 10^{-2}$.

\section{Reference metric \texorpdfstring{$\GQbar$}{Gbar}}
\label{sec:reference-metric}

Constant symmetric positive-definite matrix defining the OU stationary covariance $\GQbar^{-1}$. \textbf{Independent of $u$, $z_e$, $c$, and sample.}

\paragraph{Three options.}
\begin{itemize}[noitemsep]
    \item $\GQbar = I_{n_q}$: simplest, strongest data-independence claim. Stationary $\N(0, I_{n_q})$ isotropic.
    \item $\GQbar = \GQ(0, \bar z_e)$ for fixed $\bar z_e$: tracks geometry at chart center, introduces hyperparameter $\bar z_e$.
    \item $\GQbar = \E_{u_0 \sim p_\text{data}, z_e}[\GQ(u_0, z_e)]$: best forward-marginal match, but data-statistic-dependent (sampling-time still constant matrix, computed once).
\end{itemize}
Default: $\GQbar = I_{n_q}$. Ablate other options.

\section{Trajectory product manifold}
\label{sec:product-manifold}

$\tau = (x_0, \ldots, x_H)$, $u_{0:H} \in \R^{n_q(H+1)}$. Product reference: $\bar G_\text{traj} = \mathrm{diag}(\GQbar, \ldots, \GQbar)$.

\section{Forward SDE: variance-preserving OU on chart}
\label{sec:forward-sde}

\paragraph{Chart-coordinate forward SDE.}
\begin{equation}
\boxed{ du_r = -\tfrac{1}{2}\beta(r)\, u_r\, dr + \sqrt{\beta(r)}\, \GQbar^{-1/2}\, dW_r, \quad u_0 \sim p_\text{data}, \quad r \in [0, K]. }
\end{equation}

\subsection{Positioning: chart-space OU, not Riemannian Brownian on \texorpdfstring{$\Mbpose$}{M}}
\label{sec:forward-sde-positioning}

\paragraph{Honest characterization.} The forward SDE is a constant-coefficient OU process on the chart $u \in \R^{n_q}$. The diffusion matrix $\GQbar^{-1/2}$ is constant (independent of $u$); there is no Christoffel correction, no Laplace-Beltrami term, no $u$-dependent metric in the noise. This is \emph{not} a metric-aware Riemannian Brownian motion under $\GQ(u, z_e)$.

\paragraph{Why this is the design choice, not a defect.} A true Riemannian OU under $\GQ$ would lose the closed-form Gaussian transition kernel; the small-$r$ behavior would revert to Varadhan-style approximation. The constant-coefficient choice trades intrinsic Riemannian geometry for an exact closed-form transition, which is what enables the IK-free claim. $\GQ$ still plays a role, just not as the noise covariance: it appears in loss weighting, guidance preconditioning, and tangent identities.

\paragraph{Therefore.} The system is a \emph{chart-space OU sampler on a self-model graph manifold with Riemannian-weighted loss and natural-gradient-style guidance}. The manifold $\Mbpose$ enters via graph retraction at every step, not via the forward noise covariance.

\subsection{Closed-form transition kernel}

For constant diffusion matrix, the OU SDE admits a Gaussian transition:
\begin{equation}
\boxed{ p_{r \mid 0}(u_r \mid u_0) = \N\!\bigl(u_r;\, \alpha(r)\, u_0,\, \sigma^2(r)\, \GQbar^{-1}\bigr), }
\end{equation}
\begin{equation}
\alpha(r) := e^{-\tau(r)/2}, \qquad \sigma^2(r) := 1 - e^{-\tau(r)}, \qquad \tau(r) := \int_0^r \beta(s)\, ds.
\end{equation}

\paragraph{Exact Euclidean score.}
\begin{equation}
\boxed{ \nabla_{u_r} \log p_{r \mid 0}(u_r \mid u_0) = -\frac{\GQbar\, (u_r - \alpha(r)\, u_0)}{\sigma^2(r)}. }
\end{equation}
Closed-form, exact at all $r$, no Varadhan approximation.

\paragraph{Stationary distribution.}
\begin{equation}
p_\infty(u) = \lim_{r \to \infty} p_r(u) = \N(0, \GQbar^{-1}).
\end{equation}
Closed-form, data-independent, conditioning-independent.

\subsection{$\beta$-schedule and finite-$K$ mismatch}
\label{sec:beta-schedule}

Linear schedule: $\beta(r) = \beta_0 + r(\beta_f - \beta_0)$, $\beta_0 = 10^{-3}$, $K = 1$.

Cumulative: $\tau(K) = \beta_0 K + \tfrac{1}{2}(\beta_f - \beta_0) K^2$.

\paragraph{Corrected reference table.}

\begin{center}
\begin{tabular}{r|r|r|l}
$\beta_f$ & $\tau(K)$ & $\alpha(K) = e^{-\tau(K)/2}$ & interpretation \\
\hline
4 & $\approx 2.00$ & $\approx 0.368$ & terminal mismatch large \\
10 & $\approx 5.00$ & $\approx 0.082$ & substantial reduction \\
20 & $\approx 10.00$ & $\approx 0.0067$ & effectively stationary \\
\end{tabular}
\end{center}

\paragraph{Important distinction.} ``OU stationary exists exactly'' is a \emph{structural} property — distinct from the finite-$K$ marginal-reference match. The former is exact at all $\beta_f$ (it is a property of the SDE, independent of how long we run it). The latter improves as $\beta_f$ increases. The IK-free claim rests on the structural property: there is a well-defined task-agnostic distribution to start the reverse SDE from. Whether that distribution exactly matches the forward marginal $p_K$ depends on the empirical $\KL(p_K \| p_\text{ref})$ at the chosen $\beta_f$.

\section{Score conventions (fixed)}
\label{sec:score-conventions}

\paragraph{Network output: Euclidean score.}
\begin{equation}
\boxed{ s_\theta^u(r, x, c, z_e) \approx \nabla_u \log p_r(u \mid c, z_e). }
\end{equation}
This is the Euclidean coordinate gradient of $\log p_r$, not the Riemannian gradient. The choice is consistent with Score-SDE (Song et al., 2021) convention.

\paragraph{Network input.}
\begin{equation}
\text{input}_h = [u_h, r, h/H, z_e, T_\text{start}, T_\text{target}].
\end{equation}
No anchor, no $q^\text{init}$, no IK-derived variable. The conditioning information $(z_e, T_\text{start}, T_\text{target})$ enters here.

\paragraph{Multimodality.} Captured by the score network alone, conditional on $c$. The v4 mode-aware IK warm-start is removed.

\section{Score matching loss}
\label{sec:loss}

\paragraph{Chart-block target (exact Euclidean OU score).}
\begin{equation}
\boxed{ s^{u, *}(u_r, u_0; r) = -\frac{\GQbar\, (u_r - \alpha(r)\, u_0)}{\sigma^2(r)}. }
\end{equation}
Note: \emph{no} $\GQ^{-1}$ factor (v5 error corrected). This is exactly $\nabla_{u_r} \log p_{r \mid 0}(u_r \mid u_0)$ for the closed-form OU transition.

\paragraph{Primary loss: closed-form OU score matching.}
\begin{equation}
\boxed{ \Loss_\text{score} = \E_{r, \tau_0, \tau_r}\!\left[\, w(r) \sum_h \bigl\| s_{\theta, h}^u - s^{u, *}_h(u_{h, r}, u_{h, 0}; r) \bigr\|_{M(u_{h, r}, z_e)}^2 \,\right]. }
\end{equation}
where $M$ is a weighting matrix. Choices for ablation:
\begin{itemize}[noitemsep]
    \item $M = I_{n_q}$: standard Euclidean.
    \item $M = \GQ^{-1}$: down-weights metric-anisotropic directions.
    \item $M = \GQ$: up-weights metric-anisotropic directions.
\end{itemize}
\textbf{The target $s^{u, *}$ does not change across these choices; only the weighting metric of the residual.} Default: $M = I_{n_q}$, with $w(r) = \sigma^2(r)^2$ (SNR-aware).

\paragraph{Auxiliary pose-geometric consistency regularizer (optional).} The pose component $T = \widetilde T_\phi(u, z_e)$ is a deterministic function of $u$; therefore the $u$-space OU does \emph{not} induce a closed-form Gaussian distribution on $T$, and the $\Log_{\SE(3)}$ term is not an exact transition score. We retain it as a Varadhan-style consistency regularizer:
\begin{equation}
\Loss_\text{pose} = \E_{r, \tau_0, \tau_r}\!\left[\, \lambda_p(r) \sum_h \Bigl\| \JQ(u_{h, r}, z_e)\, s_{\theta, h}^u - \tfrac{1}{\tau(r)} \Log_{\SE(3)}\!\bigl(\widetilde T_\phi(u_{h, r})^{-1} \widetilde T_\phi(u_{h, 0})\bigr) \Bigr\|_W^2 \,\right].
\end{equation}
This term enforces that the lifted score $\JQ s_\theta^u$ aligns with the pose-tangent displacement implied by demos. It is an auxiliary geometric prior, not an exact DSM target.

\paragraph{Total loss.}
\begin{equation}
\Loss = \Loss_\text{score} + \mu_\text{pose}\, \Loss_\text{pose}, \quad \mu_\text{pose} \in [0, 1].
\end{equation}
$\mu_\text{pose} = 0$ gives clean ``exact closed-form OU score matching'' baseline; $\mu_\text{pose} > 0$ adds pose consistency. \textbf{Ablate.}

\paragraph{Validity regime of $\Loss_\text{pose}$.} The Varadhan-style pose term is accurate for small $r$ (small $u_0$-$u_r$ displacement) and degrades at large $r$. The $\lambda_p(r)$ schedule should down-weight large-$r$ contributions; default $\lambda_p(r) = \mathbb{1}[\tau(r) < \tau_\text{cutoff}]$ with $\tau_\text{cutoff} \approx 0.5$.

\section{Sampling}
\label{sec:sampling}

\paragraph{Sampler positioning.} Chart-space OU sampler with graph retraction. \textbf{Not} an intrinsic Riemannian Brownian motion (the term ``geodesic random walk'' from v5 is dropped). Graph retraction onto $\Mbpose$ ensures learned-manifold consistency (see \S\ref{sec:feasibility-tiers} for what this means and does not mean).

\paragraph{Graph retraction.}
\begin{equation}
\Retr_{(u, T)}(\delta u) = \bigl(u + \delta u,\, \widetilde T_\phi(u + \delta u, z_e)\bigr).
\end{equation}

\subsection{Reference distribution}
\label{sec:reference-distribution}

\begin{equation}
\boxed{ p_\text{ref}(\tau) = \prod_{h=0}^H \N(u_h;\, 0,\, \GQbar^{-1}) = \prod_h p_\infty. }
\end{equation}

\paragraph{Justification.} $p_\text{ref} = p_\infty$ by definition. At finite $K$, $p_K \neq p_\text{ref}$ exactly; the gap is measured by $\KL(p_K \| p_\text{ref})$ or $W_2$. The IK-free claim rests on $p_\text{ref}$ being data-independent and conditioning-independent at sampling time — not on $p_K = p_\text{ref}$ at the chosen $\beta_f$. The latter is an empirical question, addressed by the $\beta_f$ sweep (\S\ref{sec:beta-schedule}).

\paragraph{Sampling-time independence audit.}
\begin{itemize}[noitemsep]
    \item Does $p_\text{ref}$ depend on $z_e$? \textbf{No} (when $\GQbar = I_{n_q}$; depends on $\bar z_e$ but not sample-time $z_e$ for $\GQbar = \GQ(0, \bar z_e)$).
    \item Does $p_\text{ref}$ depend on $T_\text{start}, T_\text{target}$? \textbf{No.}
    \item Does $p_\text{ref}$ require an IK solver call? \textbf{No.}
    \item Does $p_\text{ref}$ require $q_\text{warm}$ or mode-aware seeding? \textbf{No.}
\end{itemize}

\subsection{Reverse SDE (Convention 1: forward time \texorpdfstring{$r$}{r}, \texorpdfstring{$\Delta r < 0$}{Delta r < 0})}
\label{sec:reverse-sde}

\paragraph{Forward SDE in standard form.}
\begin{equation}
du_r = f(u_r, r)\, dr + g(r)\, L\, dW_r, \quad f(u, r) = -\tfrac{1}{2}\beta(r)\, u, \quad g(r) = \sqrt{\beta(r)}, \quad L = \GQbar^{-1/2}.
\end{equation}
$LL^\top = \GQbar^{-1}$.

\paragraph{Anderson(1982) reverse SDE in forward time $r$.}
\begin{equation}
\boxed{ du_r = \bigl[f(u_r, r) - g^2(r)\, LL^\top\, \nabla_{u_r} \log p_r(u_r)\bigr]\, dr + g(r)\, L\, d\bar W_r, \quad dr < 0. }
\end{equation}

Substituting:
\begin{equation}
du_r = \bigl[-\tfrac{1}{2}\beta(r)\, u_r - \beta(r)\, \GQbar^{-1}\, s_\theta^u(u_r, r, c, z_e)\bigr]\, dr + \sqrt{\beta(r)}\, \GQbar^{-1/2}\, d\bar W_r, \quad dr < 0.
\end{equation}

\paragraph{Discrete update (Euler-Maruyama, Convention 1).}
\begin{equation}
\boxed{ u_{k+1} = u_k + \bigl[-\tfrac{1}{2}\beta(r_k)\, u_k - \beta(r_k)\, \GQbar^{-1}\, s_\theta^u(u_k, r_k, c, z_e)\bigr]\, \Delta r + \sqrt{\beta(r_k) |\Delta r|}\, \GQbar^{-1/2}\, \xi_k, }
\end{equation}
where $\Delta r < 0$, $\xi_k \sim \N(0, I_{n_q})$, $r_{k+1} = r_k + \Delta r < r_k$. The integration starts at $r_0 = K$ with $u_0 \sim p_\text{ref}$ and proceeds toward $r \to 0$.

\paragraph{Sign verification.} At each step, $-\tfrac{1}{2}\beta u\, \Delta r$ with $\Delta r < 0$ gives a positive contribution along $+u$ direction: this is the reverse-direction OU mirror (pushing $u$ away from origin, mirroring the forward attraction). The score term $-\beta \GQbar^{-1} s_\theta^u \cdot \Delta r$ with $\Delta r < 0$ gives a positive contribution along $+\nabla \log p$ direction: this points toward higher-density regions, as required. Signs are consistent.

\subsection{Reverse step with multi-component guidance}
\label{sec:sampling-step}

Adding guidance as preconditioned additive correction to the score:
\begin{equation}
\mu_k^u = -\tfrac{1}{2} u_k - \GQbar^{-1}\, \bigl[s_\theta^u(u_k, r_k, c, z_e) + m_k \cdot \GQ(u_k, z_e)^{-1}\, \nabla_u R_\text{total}(\tau_k, c)\bigr]
\end{equation}
\begin{equation}
\delta u_k = \beta(r_k)\, \mu_k^u\, \Delta r + \sqrt{\beta(r_k) |\Delta r|}\, \GQbar^{-1/2}\, \xi_k
\end{equation}
\begin{equation}
u_{k+1} = u_k + \delta u_k, \quad x_{k+1} = \Hbpose(u_{k+1}, z_e).
\end{equation}

\paragraph{Status of $\GQ^{-1}$ in guidance.} The factor $\GQ^{-1}$ in front of $\nabla R$ is a natural-gradient-style preconditioning — it scales the reward gradient by the inverse induced metric. \textbf{This is a heuristic, not from the exact reverse SDE.} It is a deliberate design choice that respects the manifold geometry; alternative choices ($M = I$, $M = \GQbar^{-1}$) are valid for ablation. The score correction term $-\GQbar^{-1} s_\theta^u$ does come from the exact reverse SDE.

\paragraph{Three roles, three terms.}
\begin{enumerate}[noitemsep]
    \item Reverse OU mirror $-\tfrac{1}{2} u_k$: data-independent, defines the diffusion process.
    \item Score correction $-\GQbar^{-1} s_\theta^u$: data-conditional, reconstructs $p_\text{data}$.
    \item Guidance $-\GQbar^{-1} \GQ^{-1} \nabla R$: task-specific, injects $T_\text{start}, T_\text{target}$ via reward.
\end{enumerate}

\section{Multi-component guidance}
\label{sec:guidance}

\paragraph{Status as conditioning channel.} The score network conditioning $(T_\text{start}, T_\text{target}, z_e)$ is the \emph{primary} channel by which task information enters sampling. Reward guidance is an \emph{additional, non-IK, non-network explicit mechanism}. It is not the sole channel.

\paragraph{Trajectory-level guided update.}
\begin{equation}
s_\text{guided}^u = s_\theta^u + G_\text{traj}^{-1}\, \nabla_{u_{0:H}} R_\text{total}(\tau, c),
\end{equation}
where $G_\text{traj} = \mathrm{diag}(\GQ(u_h, z_e))$. As noted, $G_\text{traj}^{-1}$ preconditioning is heuristic.

\paragraph{Reward.}
\begin{equation}
R_\text{total} = \alpha_s R_\text{start} + \alpha_g R_\text{goal} + \alpha_v R_\text{vel} + \alpha_a R_\text{acc}.
\end{equation}

\paragraph{Pose start anchor.}
\begin{equation}
e_\text{start}(u_0, z_e) = \Log_{\SE(3)}\bigl(\widetilde T_\phi(u_0, z_e)^{-1} T_\text{start}\bigr), \quad R_\text{start} = -(\alpha_p^s \|e_\rho^s\|^2 + \alpha_R^s \|e_\omega^s\|^2).
\end{equation}
Recommended: $\alpha_s \geq 2 \alpha_g$ (start anchor compensating for absence of IK warm-start).

\paragraph{Goal, velocity, acceleration.} As in v4/v5.

\paragraph{Mask schedules.} $m_h^s = \mathbb{1}[h \leq H/4]$, $m_h^g = \mathbb{1}[h \geq 3H/4]$, $m_h^v = m_h^a = 1$.

\section{Feasibility claims, tiered}
\label{sec:feasibility-tiers}

\paragraph{Tier 1: Joint feasibility — exact by construction.}
\begin{equation}
\psi(u) \in (q_\text{min}, q_\text{max}) \quad \forall u \in \R^{n_q}.
\end{equation}
This holds for any sample, any training distribution, any extrapolation. Structurally guaranteed by the chart map.

\paragraph{Tier 2: Learned-manifold adherence — exact w.r.t.\ $\Tphi$.}
\begin{equation}
\|g_\phi(x, z_e)\|_W = \|\Log_{\SE(3)}(\Tphi(\psi(u), z_e)^{-1} T)\|_W = 0
\end{equation}
holds for samples produced by $\Hbpose$ retraction. This says $T$ is consistent with $\Tphi$ at $q = \psi(u)$. It does \emph{not} say $T$ matches the true robot pose.

\paragraph{Tier 3: True physical feasibility — depends on self-model accuracy.} The true robot end-effector at configuration $q$ has pose $T_\text{true}(q, z_e)$, which may differ from $\Tphi(q, z_e)$ by the self-model approximation error. Stage-1 trained $\xi_\phi$ residual on Franka achieves $0.003$ m / $0.002$ rad inside-training and $0.004$ m / $0.003$ rad outside-1$\times$ (REPORT \S3). Tier 3 feasibility is bounded by these errors, plus dynamics/torque/collision considerations not modeled here.

\paragraph{What v5.1 does and does not guarantee.}
\begin{itemize}[noitemsep]
    \item Guarantees joint limit satisfaction (Tier 1).
    \item Guarantees consistency with learned self-model (Tier 2).
    \item Does \emph{not} guarantee collision-free trajectories.
    \item Does \emph{not} guarantee dynamic feasibility (torque, velocity, acceleration limits beyond the soft chart-smoothness reward).
    \item Does \emph{not} guarantee that $T$ matches the true robot pose (bounded by Stage-1 self-model error).
\end{itemize}

\section{Evaluation metrics}
\label{sec:eval}

Standard metrics from v4/v5: $e_p, e_R, \text{succ}@(\rho_p, \rho_R)$, mode capture, smoothness.

\paragraph{Required new metrics for v5.1.}
\begin{itemize}[noitemsep]
    \item \textbf{Reference-marginal match.} $\KL(p_K \| p_\text{ref})$ or $W_2(\text{empirical}(u_K), \N(0, \GQbar^{-1}))$. This validates the IK-free reference choice for the chosen $\beta_f$. Required because the structural exactness of $p_\infty$ does not imply finite-$K$ exactness of $p_K$.
    \item \textbf{Joint limit violation rate.} Should be exactly $0$ (Tier 1 verification, sanity check).
    \item \textbf{Chart saturation $\|u\|_\infty^{p99}$.} Healthy: $\lesssim 2$. Saturation: $> 3$.
    \item \textbf{Tier 3 self-model error proxy.} Difference between $\widetilde T_\phi(u_H, z_e)$ and a high-fidelity reference pose, when available; otherwise reported as ``not measured.''
\end{itemize}

\section{Pipeline summary}

\begin{align}
\text{Stage 1.}\quad & \Tphi = T_\text{analytic} \cdot \exp_{\SE(3)}(\xi_\phi^\wedge) \\
\text{Stage 2.}\quad & \Mbpose(z_e), \quad \GQ = I + {\JQ}^\top W \JQ, \quad \GQbar = I_{n_q} \\
\text{Stage 3.}\quad & u_h^{(i)} = \psi^{-1}(q_h^{(i)}) \text{ (demo projection)} \\
\text{Stage 4.}\quad & du_r = -\tfrac{1}{2}\beta\, u_r\, dr + \sqrt{\beta}\, \GQbar^{-1/2}\, dW_r \\
\text{Stage 5.}\quad & p_{r \mid 0} = \N(\alpha(r) u_0, \sigma^2(r) \GQbar^{-1}) \text{ (closed-form)} \\
\text{Stage 6.}\quad & \min_\theta \Loss_\text{score} + \mu_\text{pose} \Loss_\text{pose} \\
\text{Stage 7.}\quad & u^{(K)} \sim \N(0, \GQbar^{-1}) \text{ (IK-free)} \\
\text{Stage 8.}\quad & \text{Reverse SDE Convention 1: } \Delta r < 0, \text{ with score + guidance} \\
\text{Stage 9.}\quad & q_h = \psi(u_h^{(0)}), \quad T_h = \Tphi(q_h, z_e)
\end{align}

\section{Limitations}

\paragraph{Score-net capacity demand.} IK seed removal places multimodality and trajectory shape on score net alone. Higher capacity / more training steps expected vs.\ v4.

\paragraph{Pose-block term is auxiliary, not exact.} Honest framing: $\Loss_\text{pose}$ is a Varadhan-style consistency regularizer, not an exact DSM target. $\mu_\text{pose}$ should be tuned by ablation.

\paragraph{$\GQbar$ choice impacts forward-reference match.} Different $\GQbar$ choices alter how closely $p_K$ approaches $p_\text{ref}$ for given $\beta_f$. The IK-free claim is robust to $\GQbar$ choice (always data-independent at sampling), but pose-accuracy performance is not.

\paragraph{Saturation.} Chart $|u| > 3$ ill-conditioned. If observed, add $-\alpha_u \|u\|^2$ to $R_\text{total}$.

\paragraph{Interior pose accuracy expected below v4.} ABL3 in v4 showed IK seed contributed $\sim 75$ pp pose accuracy on interior demos. v5.1 cannot recover this from architectural change alone; expected boundary-active comparable or better, interior worse.

\paragraph{Tier 3 not addressed.} Collision, torque, dynamics outside framework scope.

\section{Summary of v5.1 changes (relative to v5)}

\begin{center}
\begin{tabular}{l|p{4.0cm}|p{5.5cm}}
\textbf{Issue (v5 review)} & \textbf{v5} & \textbf{v5.1 (this document)} \\
\hline
Finite-$K$ exact match & ``Exact'' claimed & Structural stationary exact; finite-$K$ $p_K \neq p_\text{ref}$, measured empirically \\
\hline
$\beta_f$-$\alpha$ table & $\beta_f = 10 \Rightarrow \alpha \approx 0.007$ (error) & $\beta_f = 10 \Rightarrow \alpha \approx 0.082$; $\beta_f = 20 \Rightarrow \alpha \approx 0.007$ \\
\hline
Score convention & Mixed (Euclidean vs.\ Riemannian) & Fixed to Euclidean: $s_\theta^u = \nabla_u \log p_r$ \\
\hline
DSM target & $-\GQ^{-1} \GQbar (u_r - \alpha u_0)/\sigma^2$ & $-\GQbar (u_r - \alpha u_0)/\sigma^2$ (no $\GQ^{-1}$) \\
\hline
Reverse SDE sign & Mixed conventions & Convention 1: forward time $r$, $\Delta r < 0$ \\
\hline
Sampler name & ``Geodesic random walk'' & ``Chart-space OU sampler with graph retraction'' \\
\hline
Pose-block term & ``Hybrid DSM target'' & ``Auxiliary pose-geometric consistency regularizer'' \\
\hline
Conditioning channel & ``Sole channel'' & ``Non-IK, non-network explicit channel'' (network is primary) \\
\hline
Feasibility claim & ``By construction'' (flat) & Tier 1 joint / Tier 2 learned-manifold / Tier 3 true physics \\
\hline
``Riemannian SDE'' claim & Implied throughout & Made precise: chart-space OU, not Riemannian Brownian \\
\hline
Reference-marginal eval & Optional & Required ($\KL$ or $W_2$) \\
\end{tabular}
\end{center}

\paragraph{What v5.1 keeps from v5.}
\begin{itemize}[noitemsep]
    \item Bounded chart with joint feasibility by construction.
    \item OU forward with closed-form Gaussian transition (now properly named ``chart-space OU,'' not Riemannian).
    \item Closed-form stationary $\N(0, \GQbar^{-1})$ enabling IK-free reference.
    \item Closed-form Euclidean OU score as the exact training target (chart block).
    \item Self-model graph manifold $\Mbpose$ via retraction.
    \item Multi-component reward guidance.
\end{itemize}

\paragraph{What v5.1 abandons from v5.}
\begin{itemize}[noitemsep]
    \item ``Exact match'' claim at finite $K$.
    \item Mixed score conventions.
    \item ``Geodesic random walk'' terminology.
    \item ``Hybrid DSM'' framing of the pose term.
    \item ``Sole conditioning channel'' framing of guidance.
    \item ``Riemannian diffusion'' overclaim.
\end{itemize}

\appendix

\section{Closed-form reward gradients on bounded chart}

\begin{equation}
\nabla_{u_H} R_\text{goal} = +2\, \JQ(u_H, z_e)^\top\, \mathcal{J}_l^{-\top}(e_\text{goal})\, W_g\, e_\text{goal}(u_H, z_e).
\end{equation}
$\JQ = \Jpose \Dpsi$ absorbs the chart chain rule. Implementation via autograd.

\section{Comparison to v4 reverse SDE}

v4: $dY = \beta s_\theta^q\, dr + \sqrt{\beta}\, \sigma\, d\bar W$, $\Delta r < 0$, no OU mirror, $\sigma\sigma^\top = G_\text{pose}^{-1}$.

v5.1: $du = [-\tfrac{1}{2}\beta u - \beta \GQbar^{-1} s_\theta^u]\, dr + \sqrt{\beta}\, \GQbar^{-1/2}\, d\bar W$, $\Delta r < 0$.

Structural addition in v5.1: OU mirror $-\tfrac{1}{2}\beta u$. This is the term that confines reverse sampling to the stationary even when score is uninformative. v4 had no such confinement; it relied on score net + IK seed.

\section{Why Convention 1}

Convention 1 ($r$ as forward time, $\Delta r < 0$ on reverse) and Convention 2 ($\bar r = K - r$, $\Delta \bar r > 0$) are equivalent by variable substitution. Convention 1 is selected because:
\begin{itemize}[noitemsep]
    \item Standard Score-SDE (Song et al., 2021) and most diffusion model implementations use it.
    \item It avoids confusion about which direction the SDE is integrated when reading the equation: $dr < 0$ is explicit.
    \item Discrete update formulas align with DDPM-style $t = T, T-1, \ldots, 1$ indexing.
\end{itemize}

\section{Future work}

\paragraph{Pose-space exact transition.} Lie-OU on $\SE(3)$ to remove the Varadhan approximation in $\Loss_\text{pose}$.

\paragraph{Anisotropic $\GQbar$.} Ablate $\GQbar \in \{I, \GQ(0, \bar z_e), \E[\GQ]\}$.

\paragraph{Comparison vs.\ v5-option-X.} Same demo, same eval, both trained from scratch. Does exact OU consistency translate to measurable improvement over Brownian + large horizon?

\paragraph{Learned reference.} $g_\psi(c, z_e) \to (\mu_\text{ref}, \Sigma_\text{ref})$ as soft trainable IK replacement.

\paragraph{Physically-grounded $W$.} Inherited from v4.

\end{document}